\documentclass[11pt]{article}
\usepackage{icmlko}
\begin{document}

\title{정확히 영이 나왔는데 이번엔 그것이 옳았다 --- ``차가 0 이면 배선을 의심한다''를 뒤집어 무효 대조로 썼다}
\author{월드모델 실험실}
\date{노트 551 · 2026년 8월 4일}
\maketitle

\begin{abstract}
노트 $548$ 이 검색 축을 통째로 끄고 도메인별로 쟀다 --- \textbf{만화
$+0.0099$($10/12$) $\cdot$ 팝업 $+0.0374$($12/12$)}. 그 둘에서 검색이
\emph{해롭다}는 읽기로 \texttt{DOMDROP} 후보를 세우고 네 팔로 갈랐다.

\begin{center}
\begin{tabular}{lrr}
\hline
 & 전역으로 끔(노트 $548$) & \textbf{그 도메인만 가림} \\
\hline
만화 & $+0.0099$ & $+0.0072$ ($10/12$) \\
\textbf{팝업} & $\mathbf{+0.0374}$ & $\mathbf{+0.0018}$ (문다) \\
\hline
\end{tabular}
\end{center}

\textbf{팝업 팔은 일곱 자가 전부 정확히 $+0.0000$($0/12$)이었다} ---
\texttt{lab/loop.py} 의 \texttt{TREND\_DROP=("팝업",)} 이 노트 $351$ 에서
이미 뺐기 때문이다(관측 $0.0\%$). \textbf{이번의 ``정확히 $0$''은 결함이
아니라 예측이었다.}
\end{abstract}

\begin{figure}[t]
\centering
\includegraphics[width=\textwidth]{figs/f1.pdf}
\caption{\textbf{가} --- 두 말이 다르다. \textbf{나} --- 예측된 $0$.
\textbf{다} --- 판은 중립, 전이는 $12/12$.}
\end{figure}

\section{``전역으로 끄면 오른다''는 ``거기서 가리면 오른다''가 아니다}

노트 $548$ 의 도메인별 표는 \emph{검색 축을 전부 없앤} 판과의 차다. 그러면
\textbf{그 도메인의 예보가 바뀌는 경로가 둘}이다.

\begin{enumerate}
\item 자기 축이 없어져서
\item \textbf{남들이 검색으로 배운 것이 없어져서} --- 공유 모형이 바뀐다
\end{enumerate}

팝업은 자기 축이 애초에 없으므로 $+0.0374$ 는 \textbf{전부 ②}다. 그래서
\texttt{DOMDROP} 으로는 한 톨도 못 잡는다.

\begin{quote}
\textbf{만화도 절반쯤 그렇다} --- 전역 $+0.0099$ 인데 가림은 $+0.0072$ 다.
차 $0.0027$ 이 이웃 효과다.
\end{quote}

\section{예측된 $0$ 은 가드의 반대편이다}

노트 $509$ 가 세운 조항은 ``\textbf{차가 정확히 $0$ 이면 배선을 의심한다}''
이고, 노트 $545$ 에서 세 번째로 값을 했다. 여기서는 \textbf{같은 무늬가
사전 등록된 예측}이었다.

\begin{quote}
가드는 ``$0$ 이면 틀렸다''가 아니라 ``\textbf{$0$ 이면 이유를 대라}''이다.
이유를 미리 댈 수 있으면 그 $0$ 은 \emph{장치가 제대로 돈다}는 증거가
된다 --- 일곱 자가 전부 $0/12$ 로 딱 맞은 것이 그것이다.
\end{quote}

\section{만화 팔 --- 판은 중립인데 전이가 산다}

\begin{center}
\begin{tabular}{lrrl}
\hline
자 & 전 & 후 & 차 \\
\hline
판 & $+0.4812$ & $+0.4810$ & $-0.0001$ ($5/12$, 문다) \\
날짜 통제 판 & $+0.5143$ & $+0.5152$ & $+0.0009$ (문다) \\
무리 안 판 & $+0.3961$ & $+0.3963$ & $+0.0002$ (문다) \\
\textbf{KR 만화} & $+0.6297$ & $+0.6387$ & $\mathbf{+0.0090}$ ($\mathbf{12/12}$) \\
비게임 앱 & $+0.2874$ & $+0.2885$ & $+0.0011$ (문다) \\
\hline
만화 & $+0.3313$ & $+0.3385$ & $+0.0072$ ($10/12$) \\
\hline
\end{tabular}
\end{center}

\textbf{어떤 자도 $0$ 밖 음수가 아니고}, 거부권 $446$ 은 안 걸리고,
\textbf{노트 $426$ 거울도 통과}한다(집 밖 짝 둘 다 양수). 만화 덮음이
$16.4\%$ 뿐이라 그 축의 대부분은 \emph{결측 표시자}이고, 그 표시자가
\textbf{전이를 팔고 있었다}는 읽기가 된다.

\section*{그런데 보류한다 --- 이유 셋}

\begin{enumerate}
\item \textbf{내 사전 등록 ①이 반증됐다} --- ``판이 $0$ 밖 양수''라고 적었고
$-0.0001$ 이다. 사후에 기준을 ``전이만 오르면 된다''로 바꾸면 그것이 점수
고르기다.
\item \textbf{노트 $518$} --- $0$ 근처인 자가 넷이면 \textbf{씨앗 $40$} 으로
다시 잰다.
\item \textbf{노트 $552$ 가 분모를 바꾼다} --- 웹툰 유보에서 날짜가 틀린
$61$ 행을 빼는 정정이 진행 중이라, 지금 판 수치가 곧 낡는다.
\end{enumerate}

\begin{quote}
\textbf{채택한다면 근거는 이것이다} --- 판이 안 내려가면서 집 밖 전이가
$12/12$ 로 오른다. \textbf{노트 $351$ 이 정확히 같은 모양}(판 $+0.0002$,
$7/12$ 인데 대표가 $+0.0168$, $11/12$)으로 \texttt{TREND\_DROP} 을
채택한 선례가 있다. \textbf{정정 뒤 씨앗 $40$ 으로 다시 잰다.}
\end{quote}

\section*{조항}

\begin{quote}
\textbf{``전역으로 끄면 오른다''를 ``거기서 가리면 오른다''로 읽지 않는다.}
전역 제거의 도메인별 차에는 \emph{이웃 효과}가 섞여 있다. 가림 후보는
\textbf{그 도메인에서 직접} 재야 하고, 자기 축 덮음이 $0\%$ 인 도메인은
애초에 후보가 아니다(노트 $525$ 의 조건 ①과 같은 말).
\end{quote}

\section*{남은 것}

\begin{itemize}
\item \textbf{팝업의 $+0.0374$ 는 어떻게 잡나} --- 이웃 효과라
\texttt{DOMDROP} 밖이다. 남들의 검색 학습이 팝업 예보에 안 닿게 하려면
도메인별 모형이나 가중이 필요한데, 그것은 다른 손잡이다.
\item \textbf{만화 덮음 $16.4\%$ 의 정체} --- 원천 파일 기준 $0.8\%$ 인데
하네스에서 $16.4\%$ 다. \texttt{zero\_is\_data} 가 ``묻고 $0$ 을 받았다''를
관측으로 세기 때문인데, \emph{그 구분이 무엇을 사는지} 따로 안 봤다.
\end{itemize}

\end{document}
